\part{The Experimental Program}

The chapters of this part are protocols. Each states a hypothesis, the construction of the data, the systems compared, the measurements, the controls, and the criteria under which the hypothesis is counted as supported or as refuted. No experimental outcome is reported in this monograph. The algebraic checks of Appendix~\ref{app:pilot} verify identities and bounds of Parts~II and III and are not evaluations of any claim made here. Where a protocol uses a quantity from the literature, it is quoted as reported by its source.

Two rules apply throughout. The first is the matched-baseline rule of Theorem~\ref{thm:matched}: every comparison with a complex-valued system includes a real-valued system of equal information content, because on the sign-phase subfamily the resonance score reduces to a norm-weighted cosine, and any advantage must therefore be located in the norm factor, in the operators and feature map, or in learning. The second is a requirement of pre-specification: the primary metric, the analysis, and the acceptance threshold are fixed before data are generated, and results for which they were chosen after inspection are labeled exploratory.

Sample sizes are determined by the comparison of two proportions. For a binary success outcome with success probabilities $p_0$ and $p_1$ under two independent arms, a two-sided test at level $\alpha$ with power $1-\beta_{\mathrm{p}}$ requires approximately
\begin{equation}\label{eq:samplesize}
n\ \ge\ \frac{(z_{1-\alpha/2}+z_{1-\beta_{\mathrm p}})^2\bigl(p_0(1-p_0)+p_1(1-p_1)\bigr)}{(p_1-p_0)^2}
\end{equation}
queries per arm, and the paired design used when the same queries are posed to both systems reduces this by the correlation between the arms' outcomes. These are planning formulas and the actual effect sizes are unknown until the experiments are run.

\chapter{Reproduction of the Wave-Based Retrieval Paper}\label{ch:repro}

The first experiment reproduces the reported behavior of Listopad's system \cite{listopad2025} before any extension is tried. The paper represents patterns as complex waveforms $\psi(x)=A(x)e^{\iu\phi(x)}$, scores pairs by the resonance score of Chapter~\ref{ch:resonance}, defines the semantic operators NEG, SHIFT+, INT\_UP and INT\_DOWN on a compact vocabulary, and reports distance histograms, a heatmap of pairwise distances, and retrieval latency. The operators are not given formal definitions in the paper; the candidate definitions of Definition~\ref{def:cand-ops} are therefore the object of the reproduction, and the first deliverable is a report of whether the candidates reproduce the qualitative structure the paper describes, namely that cosine distances for the operator outputs concentrate near a common value while resonance distances separate them into blocks, with NEG aligned with the anti-phased opposite and INT\_UP and INT\_DOWN separated by amplitude.

The latency measurements are reproduced under the reported conditions: pattern length $L=1024$, a store of 500{,}000 patterns, a scalar Java backend, twelve threads, and cold cache. The paper reports an average top-one latency of approximately 135 ms with a 95th percentile near 152 ms, and comparable figures for top-ten and top-hundred queries. The reproduction records the same statistics on its own hardware, with the hardware and software configuration reported alongside. Because the paper describes latency as growing linearly with the dataset size while the reported measurement is at one size, the reproduction measures at several store sizes, from $10^3$ to $5\times10^5$, and fits the dependence; linear growth is expected for an exhaustive scan and is a hypothesis to be tested, not an assumption.

The reproduction includes the comparison the original does not make. Against the unmodified real-valued embeddings, the proposal requires matched real-valued representations of equal information capacity. Concretely, for each complex pattern in $\C^n$ the matched baseline uses a real vector in $\R^{2n}$ with the same number of free parameters, and retrieval by the norm-weighted cosine $\rho\,\cos$ and by the plain cosine. By Theorem~\ref{thm:realif} a retrieval operation that is complex-linear is a special case of a real-linear one, so the baseline is at least as expressive. A reproduction that finds the complex system superior to unmodified cosine retrieval and not superior to the matched baseline has localized the effect to the norm factor, the operators, or the feature map, which is itself a result of the program.

\paragraph{Acceptance.} The reproduction succeeds if the qualitative structure of the operator results is recovered under the candidate definitions and the latency is within a stated factor of the reported figure on comparable hardware. It does not succeed or fail on the question of superiority, which belongs to the later experiments.

\chapter{Ecphoric Pattern Completion}\label{ch:completion}

The hypothesis is that Mem8 recovers the correct episode, not merely the nearest stored vector, when the cue is partial, observations are missing, temporal context is corrupted, or competing historical explanations exist. The experiment is designed so that nearest-neighbor retrieval has a determinate failure mode, which the sufficiency theorem identifies: an episode set is constructed in which several episodes share their content and differ in order, context or outcome, so that retrieval by content alone cannot succeed for the target recollection of Definition~\ref{def:suff}.

Episodes are generated from a controlled grammar with $V$ items, episodes of length $\ell$, and a context label drawn from a finite set. Cues are formed by deleting a fraction of the episode, replacing items by distractors, shifting the temporal code, or presenting an ambiguous partial sequence compatible with two or more stored episodes. The target recollection function $T(h,c)$ is specified for each cue type, and the encoding is audited with the sufficiency criterion before any decoder is run: if the criterion fails for the target, the experiment reports the identifiability bound of Proposition~\ref{prop:identlimit} as the ceiling, and a decoder reaching it is credited with optimal performance, while a decoder above it is a bug. Success is measured by the exact-match rate of the reconstructed episode, by the distance of the reconstruction to the target in the realified space, and by the rate at which the returned episode is a stored episode that differs from the target, which separates lure responses from noise. Systems compared are nearest-neighbor retrieval with cosine, with resonance, the softmax reconstruction without context, the softmax reconstruction with context and phase, and the nonoscillatory recurrent baseline. The sample size for a given contrast follows \eqref{eq:samplesize}, and cue types are analyzed separately because the lure and the noise conditions test different mechanisms.

\chapter{Temporal Sequence Reconstruction}\label{ch:sequence}

The hypothesis concerns the recovery of event order, intervals, nested episodes and temporal boundaries, and it is tested for the claim of Chapter~\ref{ch:binding} that a multi-frequency phase code supports exact recovery of lag within the period of the code. For frequencies $\omega_1,\dots,\omega_J$ with integer periods $T_j$, the lag $t$ is decodable modulo $\operatorname{lcm}(T_j)$ by the Chinese remainder theorem when the periods are pairwise coprime, and the experiment verifies the decoding limit and measures how it degrades with phase noise.

The data are sequences of items with specified onset times, nested structures in which one episode contains others, and boundaries marked by context shifts. Tasks include order judgment between two items, interval estimation, retrieval of the item at a given lag from a cue, and segmentation of a continuous stream into episodes. Systems compared are static embeddings with additive positional codes, recurrent networks, the phase-coded memory with a learned and with a fixed code, and a nonoscillatory dynamical baseline with matched state size. The controls of Chapter~\ref{ch:binding} apply: a phase-shuffled variant in which the phase code is permuted across items, which tests whether performance depends on the phase assignment, and a frozen-phase variant in which phases are fixed after encoding, which tests whether dynamics contribute beyond a static code. Because a static positional encoding is itself a phase code, a positive result for the phase-coded memory over the recurrent network is not a result for oscillation until the frozen-phase variant has been compared.

\paragraph{Segmentation as a separate test.} Temporal phase can encode order and interval but does not determine which events form a single remembered episode, so the experiment separates segmentation from sequence retrieval. Every system is run twice, once with oracle episode boundaries supplied by the generator and once with boundaries inferred by the surprise criterion of Chapter~\ref{ch:episodic} (a boundary where the prediction error of the context state exceeds an adaptive threshold). Boundary quality is scored by the F-measure of detected against true boundaries at a stated temporal tolerance. The segmentation cost is the difference in retrieval accuracy between oracle and inferred boundaries, so that a system cannot be credited for sequence retrieval that depended on boundaries it did not infer, nor charged for a segmentation error as if it were a retrieval error.

\chapter{Contextual Interference and Source Attribution}\label{ch:source-exp}

The hypothesis is that systems which store provenance can attribute source and calibrate false recollection, while systems limited to content cannot, as stated in Proposition~\ref{prop:source}. Datasets are constructed in which identical content appears in different contexts and is associated with conflicting historical outcomes. For each item a ground-truth set of sources is recorded. Queries ask for the content, for the source, and for whether the item was previously encountered; the three correspond to retrieval, source discrimination and familiarity.

The measurements are source accuracy, the rate of false recollection (a confident report of an item or source that was not encountered), calibration of reported confidence measured by the expected calibration error and by the reliability diagram, and the sensitivity of each to changes of context between encoding and retrieval, using the encoding-specificity prediction $S=(1+\cos(z-z_0))/2$ of Proposition~\ref{prop:encspec}. The centroid lure of Proposition~\ref{prop:lure} is tested directly: stores are constructed with a cluster of related items and a lure at the cluster mean, and the rate at which the system endorses the lure as previously seen is recorded as a function of the inverse temperature $\beta$ and the number of related items. The systems compared are a content-only softmax memory, the same with context, the same with the ledger of Chapter~\ref{ch:provenance}, and a baseline that stores explicit source labels, which sets the ceiling for source accuracy. The prediction following from the sufficiency theorem is that source accuracy of any system without access to source information is bounded by the identifiability limit, and the experiment treats any exceedance of that bound as evidence of leakage in the dataset or the pipeline.

\paragraph{Warrant coverage.} In addition to source accuracy, the experiment records the warrant coverage $\omega$ of Definition~\ref{def:warrant} for each reconstruction, and the accuracy of reported confidence separately for warranted and inferred details. A system whose details are mostly inferred but reported with the confidence of warranted ones is miscalibrated in a way that source accuracy on the final answer alone does not reveal.

\chapter{Consolidation and Reconsolidation}\label{ch:reconsol-exp}

The hypothesis is that evidence-gated updating allows adaptation to new evidence while keeping the original record recoverable and resisting unwarranted revision, as argued in Chapter~\ref{ch:reconsol}. Streams of episodes are generated with a latent state that changes at known times, so that earlier interpretations become obsolete. Three types of events are interleaved: confirming evidence, contradicting evidence of controlled strength, and reconstructions retrieved by the system itself and fed back as cues without external evidence.

Two quantities are measured and reported together. Adaptability is the speed with which the system's reconstruction follows a change in the world after contradicting evidence, measured by the number of events to reach a stated accuracy on the new state. Resistance to unwarranted revision is the distortion of stored traces after a sequence of self-generated retrievals with no external evidence, measured as the distance from the original inscription; by Proposition~\ref{prop:drift}, an update rule driven by the system's own output collapses traces toward the reconstruction, and the experiment tests whether the evidence gate prevents it. The recoverability of the original record is tested by reversal: the earlier state is reconstructed from the current state and the ledger, and compared bitwise with the earlier state (Proposition~\ref{prop:reverse}). The systems compared are ungated self-referential updating, evidence-gated updating, and a baseline that never updates, which sets the floor for adaptability and the ceiling for resistance. The result of interest is the trade-off curve between adaptability and resistance as the gate threshold varies, and the claim of the monograph is that the gated system dominates the ungated one on the resistance axis at matched adaptability.

\chapter{Oscillatory Ablation Experiments}\label{ch:ablation}

This is the decisive experiment for the necessity of oscillatory coordination. The question, as posed in the proposal, is whether coordinated oscillatory dynamics improve reconstruction under matched parameter counts, computational budgets and information content. The experiment is designed so that a negative answer is as informative as a positive one.

The full system is compared against five variants that differ from it in exactly one module of Chapter~\ref{ch:oscengine}. The phase-frozen variant fixes phases at their initial values. The phase-randomized variant draws phases from the same marginal distribution independently of state. The frequency-uncoupled variant sets $K=0$ so that each oscillator runs at its own frequency. The nonoscillatory recurrent variant replaces the Kuramoto dynamics by a learned map with the same parameter count and no sinusoidal structure. The static resonance variant removes the contextual operator altogether, so that scoring uses the content alone. Tasks are those of Chapters~\ref{ch:completion} to~\ref{ch:source-exp}, which cover completion, temporal order and source attribution, and the primary outcome is reconstruction accuracy on the target recollection function under each task's cue family.

Matching is the delicate part and is specified in advance. Parameter counts are matched by adjusting the hidden width of the nonoscillatory variant. Computational budget is matched by the number of floating-point operations per retrieval, counted analytically. Information content is matched by holding the encoding $E$ fixed across variants, so that Theorem~\ref{thm:suff} gives the same ceiling to all of them and differences are attributable to the decoder. The analysis uses paired comparisons across a common set of queries with a mixed-effects model for task and seed, and the acceptance rule is stated before the run: oscillation is judged to contribute if the full system exceeds every one of the five variants by at least a pre-specified margin on at least a pre-specified number of tasks, with the seeds as the unit of replication. A result in which the nonoscillatory recurrent variant matches the full system is reported as showing that the oscillatory realization of the computation is not necessary for it, which is consistent with the position of Chapter~\ref{ch:equivalence} that the necessary properties are those of the computation and not its physical substrate.

\paragraph{An integrated longitudinal benchmark.} The component experiments leave open whether the architecture preserves reconstructive capacity over time or merely contains several independently functional subsystems, so the program includes one benchmark that requires the capacities together. A system encounters partially overlapping episodes, later receives misleading information about one of them, undergoes consolidation and compression, and is then asked to reconstruct a particular episode from an incomplete cue. It must recover temporal order, identify the correct historical source, distinguish original evidence from subsequent interpretation, and report uncertainty. The cycle of misleading information, update, compression and query is repeated, and every metric above is reported as a function of the number of cycles, so that the object of study is the trajectory of performance and not a single retrieval. The central comparison is against a system that has the same provenance ledger, the same recurrent inference mechanism, the same storage budget, and the same training data, and differs only in lacking the oscillatory coordination. Without that control, an improvement could be attributed to additional information or engineering capacity and not to the proposed mechanism. The prediction following from Theorem~\ref{thm:dynsuff} is that trace-only systems lose answerability irreversibly at the first compression that crosses the commitment of Corollary~\ref{cor:commit}, while ledger-backed systems recover it at the cost of replay, and the benchmark measures that cost.

\chapter{Scaling and Computational Efficiency}\label{ch:scaling}

The hypothesis is that retrieval cost and accuracy scale as predicted by the capacity analysis of Chapter~\ref{ch:capacity}. Measured quantities are capacity (the number of stored traces at which accuracy falls below a threshold), retrieval latency, throughput in queries per second, update cost per write, energy where it can be measured, and accuracy as the number of related traces grows. The predictions to be tested are specific: the expected number of false activations of Proposition~\ref{prop:falseact} as a function of the threshold, the dimension and the store size, and the exponential capacity of Theorem~\ref{thm:capacity} with its required scaling of $\beta$ with the dimension. The random-pattern model is the best case, so the experiment is repeated on structured stores with clustered traces, where overlaps are larger and the capacity predicted by the margin condition is correspondingly smaller.

Both exact and approximate retrieval are measured. Exact retrieval is the exhaustive scan, with latency linear in the store size at fixed dimension. Approximate retrieval uses an index to restrict activation to a candidate set, and the experiment reports the recall of the true top candidate as a function of the candidate size, because an approximate method that discards the correct trace before refinement cannot be rescued by the refinement. Implementations are compared on CPU and GPU, with memory-mapped storage for stores larger than memory, and latency is reported with percentiles and under cold and warm cache separately. Interference is varied by adding competitor traces and measuring the decline of the target weight against the dilution formula of Proposition~\ref{prop:dilution}. Efficiency claims are reported with the hardware and with the cost of the baseline, since the comparison of latency across implementations on different hardware is not informative.

\chapter{Robustness and Failure Analysis}\label{ch:failure}

The program requires explicit reporting of failure. Successful resonance in a favorable condition is not evidence of general mnemonic competence, and the experiment is organized around the conditions under which the system is expected to fail, so that the failures are located and characterized.

Seven conditions are examined. Catastrophic interference is tested by the sequential learning of tasks with and without replay, measuring the accuracy on earlier tasks. False attractors are mapped by counting fixed points of the retrieval iteration that are not within a stated distance of any stored trace, as a function of $\beta$, $P$ and the overlap structure, against the stability criterion of Theorem~\ref{thm:contract}. Unstable synchronization is examined by driving the oscillator toward global synchrony and measuring the loss of selectivity. Phase noise is added at increasing levels to the phase code, and the degradation is compared with the analytic dependence of the score on phase mismatch. Adversarial cues are constructed by optimizing a perturbation of bounded norm to maximize the weight of a chosen wrong trace; since the weights depend on scores through the Lipschitz map of Proposition~\ref{prop:distort}, the bound gives the minimum perturbation needed to cross a margin, and the experiment measures how far the empirical attack is from that bound. Provenance corruption is tested by removing or altering ledger records and measuring the system's ability to detect the inconsistency, through the check that a replay reproduces the live state. Overconfident reconstruction is measured as the gap between reported confidence and accuracy on cues for which the target is not recoverable by Theorem~\ref{thm:suff}; a system that reports high confidence there is miscalibrated with respect to its own encoding. Each failure is reported with the conditions, the frequency, and the proposed mechanism, and unexplained failures are listed as such.
